\documentclass[a4paper,12pt]{report}

\usepackage[lmargin=2.00000cm,rmargin=2.0000cm,tmargin=5.5000cm,bmargin=2.500000cm,headheight=3.5cm]{geometry}        %Flexible and complete interface to document dimensions
\usepackage[utf8]{inputenc}
\usepackage[T1]{fontenc}
%\usepackage[latin1]{inputenc}
\usepackage{amsmath}
\usepackage{amsfonts}
\usepackage{amssymb}
\usepackage{lmodern}
\usepackage{float}
\usepackage{graphicx}

\usepackage[english,french]{babel}        %use for the below package `datetime'
\usepackage[babel=true]{csquotes}
\usepackage{datetime}        %Change format of `\today' with commands for current time
\renewcommand{\dateseparator}{-}
\newcommand{\headertoday}{\twodigit\day \dateseparator \twodigit\month  \dateseparator \the\year}

\usepackage{animate}
\usepackage{lastpage}
\usepackage{array}
\usepackage{lastpage}
\usepackage{multirow}
\usepackage{titling}
\usepackage{placeins}
\usepackage{media9}
\usepackage{eurosym}
\usepackage{pdflscape}
\usepackage{color}
\usepackage[table]{xcolor}
\definecolor{mon_bleu}{rgb}{0.137,0.466,0.741}
\definecolor{mon_vert}{rgb}{0.07843,0.4627,0.07843}
\definecolor{mon_rouge}{rgb}{0.62745,0.16078,0.27058}		
%\usepackage{subfigure}
\usepackage{subcaption}

\usepackage[ 
           hidelinks,
           colorlinks=true,
           linkcolor=blue,          % color of internal links (change box color with linkbordercolor)
    	   citecolor=blue,        % color of links to bibliography
    	   filecolor=blue,      % color of file links
    	   urlcolor=blue,           % color of external links
           pdfhighlight =/O]{hyperref}																		% dvipdfm package pour lien hypertextes dans pdf (pdfhighlight--> afficher la main dans le pdf, colorlinks--> colorlinks=true afficher les liens en couleurs)
  
\usepackage{blindtext}
\usepackage[final]{pdfpages}
\usepackage[french]{cleveref}																				% à placer après hyperref attention ne pas utiliser ":" dans les labels des équations avec French babel activé et cref

\usepackage{fancyhdr}
%%%%%%%%%%%%%%%%%%%%%%%%%%%%%%%%%%%%%%%%%%%%%%%%%%%%%%%%%%%%%%%%%%%%%%%%%%
%                         FORMAT PERSONNALISES                           %
% %%%%%%%%%%%%%%%%%%%%%%%%%%%%%%%%%%%%%%%%%%%%%%%%%%%%%%%%%%%%%%%%%%%%%%%%

\parindent=0pt        %leading space for paragraphs
\pagestyle{fancy}
\renewcommand{\arraystretch}{1.5}
\renewcommand{\headrulewidth}{0pt}
\fancyhead[CE,CO,LE,LO,RE,RO]{} %% clear out all headers
\fancyhead[C]{%
\begin{tabular}{|m{3.0cm}|m{10cm}|c@{}|}
\hline
\multirow{2}{*}{\includegraphics[scale=0.045]																			% LOGO LABO
{Logo-Symme.png}}& \centering \multirow{2}{*}{ \Large{\thetitle}}  &  \multirow{2}{*}{Page: \thepage ~/ \pageref{LastPage}}\\
& &  \\
\hline
%Etabli par: \newline \theauthor & \multicolumn{2}{l|}{Diffusion: interne}\\
%\hline
\multicolumn{2}{|l|}{Objet: \thetitleobject }& Date: \headertoday ~~\\ 
\hline
\end{tabular}
}
\renewcommand\footrulewidth{1pt}
\fancyfoot[L]{Document confidentiel} 																		% DOCUMENT CONFIDENTIEL
\fancyfoot[R]{Laboratoire SYMME}		 																	% NOM DU LABO

\newcolumntype{x}[1]{>{\centering\hspace{0pt}}p{#1}}
\setlength{\doublerulesep}{\arrayrulewidth} 


\usepackage{tikz}																							% pour les diagrammes
\usetikzlibrary{calc}
\usetikzlibrary{shapes.geometric,shapes.arrows,decorations.pathmorphing}
\usetikzlibrary{matrix,chains,scopes,positioning,shapes,arrows,fit}
\usetikzlibrary{calc,decorations.pathreplacing}
\usetikzlibrary{calc}
\usetikzlibrary{backgrounds,decorations.markings}

\setcounter{secnumdepth}{5}		%sous sous sous section



%%%%%%%%%%%%%%%%%%%%%%%%%%%%%%%%%%%%%%%%%%%%%%%%%%%%%%%%%%%%%%%%%%%%%%%%%%%%%%%%%%%%
%%-------------> PAGE DE GARDE INFO
%%%%%%%%%%%%%%%%%%%%%%%%%%%%%%%%%%%%%%%%%%%%%%%%%%%%%%%%%%%%%%%%%%%%%%%%%%%%%%%%%%%%

\author{Auteur}
\newcommand{\validator}{J. COLLOMB}
\title{Rapport d'avancement}
\selectlanguage{french}	
\date{\today}
\newcommand{\thetitleobject}{Comment faire des tableaux simplement ?}
\setcounter{tocdepth}{6}
\setcounter{secnumdepth}{6}


%%%%%%%%%%%%%%%%%%%%%%%%%%%%%%%%%%%%%%%%%%%%%%%%%%%%%%%%%%%%%%%%%%%%%%%%%%%%%%%%%%%%
%%-------------> DEBUT DU DOCUMENT 
%%%%%%%%%%%%%%%%%%%%%%%%%%%%%%%%%%%%%%%%%%%%%%%%%%%%%%%%%%%%%%%%%%%%%%%%%%%%%%%%%%%%

\begin{document}

\graphicspath{{Figures/}}

%%%%%%%%%%%%%%%%%%%%%%%%%%%%%%%%%%%%%%%%%%%%%%%%%%%%%%%%%%%%%%%%%%%%%%%%%%%%%%%%%%%%
%%-------------> PAGE DE GARDE 
%%%%%%%%%%%%%%%%%%%%%%%%%%%%%%%%%%%%%%%%%%%%%%%%%%%%%%%%%%%%%%%%%%%%%%%%%%%%%%%%%%%%
\begin{titlepage}
    \centering
    \vspace*{3cm}
    {\bfseries\Large
        \thetitle\\
        \thetitleobject\\
        --- \\
        Document confidentiel\\  																			% LIGNE "CONFIDENTIEL"
        \vskip2cm
    }
    \includegraphics[scale=0.65]{logos.png}  																% LOGOS DES PARTENAIRES
    \vfill
    \begin{center}
    \begin{table}[b]
    \begin{tabular}{x{.225\linewidth}|| x{.225\linewidth}|| x{.225\linewidth} || x{.225\linewidth} }
    \hline \hline
    \textbf{Date} & \textbf{Révision} & \textbf{Rédigé par} & \textbf{Validé par} \tabularnewline
    \hline \hline
    \thedate & Rev. A  &  \theauthor &  \validator \tabularnewline
    \hline
    ~ & ~ & ~ & ~ \tabularnewline
    \hline
    ~ & ~ & ~ & ~ \tabularnewline
    \hline
    ~ & ~ & ~ & ~\tabularnewline         
    \hline \hline                                                                   
    \end{tabular}
    \end{table}
    \end{center}
    \vfill
    \vfill
    \begin{tikzpicture}[remember picture,overlay]
    \node (label) at (8cm,20cm){
        \includegraphics[width=2.5cm]{Logo-Symme.png} 													% LOGO LABO
      };
    
	\draw[very thick]
		([yshift=-25pt,xshift=25pt]current page.north west)--
		([yshift=-25pt,xshift=-25pt]current page.north east)--
		([yshift=25pt,xshift=-25pt]current page.south east)--
		([yshift=25pt,xshift=25pt]current page.south west)--cycle;
    \end{tikzpicture}
\end{titlepage}


%\begin{titlepage}
%    \centering
%    
%    \vspace*{3cm}
%    {\bfseries\Large
%        \thetitle\\
%        \thetitleobject\\
%        \vskip2cm
%    }    
%    \vfill
%    \begin{center}
%    \begin{table}[b]
%    \begin{tabular}{x{.225\linewidth}|| x{.225\linewidth}|| x{.225\linewidth} || x{.225\linewidth} }
%    \hline \hline
%    \textbf{Date} & \textbf{R�vision} & \textbf{R�dig� par} & \textbf{Valid� par} \tabularnewline
%    \hline \hline
%    \thedate & Rev. A  &  \theauthor &  \validator \tabularnewline
%    \hline
%    ~ & ~ & ~ & ~ \tabularnewline
%    \hline
%    ~ & ~ & ~ & ~ \tabularnewline
%    \hline
%    ~ & ~ & ~ & ~\tabularnewline         
%    \hline \hline                                                                   
%    \end{tabular}
%    \end{table}
%    \end{center}
%    \vfill
%    \vfill
%    \begin{tikzpicture}[overlay,remember picture]
%    \node (label) at (6.7cm,19.5cm){
%        \includegraphics[width=4cm]{../Figures/logo_CT1.pdf} % also works with logo.pdf
%      };
%    
%    \draw [line width=1pt,rounded corners=7pt]
%        ($ (current page.north west) + (1.5cm,-1.5cm) $)
%        rectangle
%        ($ (current page.south east) + (-1.5cm,1.5cm) $);
%        
%%	\draw[very thick]
%%		([yshift=-25pt,xshift=25pt]current page.north west)--
%%		([yshift=-25pt,xshift=-25pt]current page.north east)--
%%		([yshift=25pt,xshift=-25pt]current page.south east)--
%%		([yshift=25pt,xshift=25pt]current page.south west)--cycle;
%    \end{tikzpicture}
%\end{titlepage}


%%%%%%%%%%%%%%%%%%%%%%%%%%%%%%%%%%%%%%%%%%%%%%%%%%%%%%%%%%%%%%%%%%%%%%%%%%%%%%%%%%%%
%%-------------> SOMMAIRE
%%%%%%%%%%%%%%%%%%%%%%%%%%%%%%%%%%%%%%%%%%%%%%%%%%%%%%%%%%%%%%%%%%%%%%%%%%%%%%%%%%%%

\renewcommand\contentsname{Sommaire}
\setcounter{chapter}{1}
\tableofcontents
%\listoffigures
%\listoftables


%%%%%%%%%%%%%%%%%%%%%%%%%%%%%%%%%%%%%%%%%%%%%%%%%%%%%%%%%%%%%%%%%%%%%%%%%%%%%%%%%%%%
%%-------------> CORPS DOCUMENT
%%%%%%%%%%%%%%%%%%%%%%%%%%%%%%%%%%%%%%%%%%%%%%%%%%%%%%%%%%%%%%%%%%%%%%%%%%%%%%%%%%%%

%-----------------------------------------------------------------------------------
%-----------------------------------------------------------------------------------
\newpage

\section{Faire des tableaux}
\subsection{Assistants des logiciels}
Certains traitements de texte, comme par exemple Texmaker, possèdent des assistants pour la génération des tableaux. La Figure~\ref{figure_assistant_tableau} est un exemple d'assistant. \\

Ces assistants permettent de créer de manière graphique le tableau, puis de générer de manière automatique le code \LaTeX associé. \\

Cependant, il est généralement nécessaire avec ces outils de devoir ajouter des "bouts de codes" pour personnaliser son tableau (allignement du texte, taille de cellule\ldots) et il est également difficile avec ces outils d'intégrer des cellules fusionnées.

\begin{figure}[hbtp]
	\centering
	\def\svgwidth{0.9\columnwidth}
	\fontsize{10pt}{10pt}\selectfont\input{Figures/assistant_tableau.pdf_tex}
	\caption{Assistant Texmaker}
	\label{figure_assistant_tableau}
\end{figure}

\subsection{Générateurs en ligne}
Des outils en ligne existent et possèdent des fonctionnalités similaires, et parfois plus développées, comme par exemple ce \href{http://www.tablesgenerator.com/latex_tables}{Générateur de tableau}. La Figure~\ref{figure_generateur_tableau} présente l'interface de cet outil en ligne. \\

\begin{figure}[hbtp]
	\centering
	\def\svgwidth{1\columnwidth}
	\fontsize{10pt}{10pt}\selectfont\input{Figures/generateur.pdf_tex}
	\caption{Générateur en ligne}
	\label{figure_generateur_tableau}
\end{figure}

Cet outil en ligne permet par exemple de :
\begin{itemize}
\item Copier/Coller des tableaux Excel/Word et de générer le code associé;
\item Générer son propre tableau;
\item Gérer la fusion des cellules;
\item Gérer les alignements;
\item Gérer les bordures;
\item \ldots
\end{itemize}


\section{Exemples}

Un exemple de tableau simple (Tableau~\ref{tableau_simple}) et un exemple de tableau complexe (Tableau~\ref{tableau_complexe}).
\begin{table}[hbtp]
\centering
\begin{tabular}{|c|c|c|}
\hline
A & B & C \\ \hline
1 & 2 & 3 \\ \hline
4 & 5 & 6 \\ \hline
\end{tabular}
\caption{Tableau simple}
\label{tableau_simple}
\end{table}


\begin{table}[hbtp]
\centering
\begin{tabular}{|c|c|c|c|c|}
\hline
\multicolumn{2}{|c|}{\multirow{2}{*}{}} & \multicolumn{3}{c|}{E} \\ \cline{3-5} 
\multicolumn{2}{|c|}{}                  & a      & b     & c     \\ \hline
\multirow{3}{*}{F}          & t         & 1      & 2     & 3     \\ \cline{2-5} 
                            & u         & 4      & 5     & 6     \\ \cline{2-5} 
                            & v         & 7      & 8     & 9     \\ \hline
\end{tabular}
\caption{Tableau complexe avec cellules fusionnées}
\label{tableau_complexe}
\end{table}

%%%%%%%%%%%%%%%%%%%%%%%%%%%%%%%%%%%%%%%%%%%%%%%%%%%%%%%%%%%%%%%%%%%%%%%%%%%%%%%%%%%%
%%-------------> FIN DU DOCUMENT
%%%%%%%%%%%%%%%%%%%%%%%%%%%%%%%%%%%%%%%%%%%%%%%%%%%%%%%%%%%%%%%%%%%%%%%%%%%%%%%%%%%%

\end{document}